%% kapa: exam=1 round=27 year=2016 period=1 form=A subject=civil_law subject_ko=민법 seq=7 number=7
%% source: Q-Net 감정평가사 기출문제 (cst00309) artlSeq=5169402 file=extracted/1008/5169402/제27회 감정평가사 1차 문제/제27회 감정평가사 1교시문제(A형).pdf
%% conversion: kordoc 4.24.51 -> tools/gen1.py (auto); review=proofread-sonnet; flags=none
\begin{question}{7}
통정허위표시의 무효를 이유로 대항할 수 없는 ‘제3자’가 될 수 있는 자를 모두 고른 것은? (다툼이 있으면 판례에 따름)
\begin{framedtext}
ㄱ. 가장의 금전소비대차에 기한 대여금채권을 가압류한 자 \par
ㄴ. 가장매매에 의한 매수인으로부터 목적 부동산에 대한 소유권이전등기 청구권 보전을 위한 가등기를 마친 제3자 \par
ㄷ. 가장매매에 의한 매수인으로부터 목적 부동산을 매수하여 소유권이전등기를 마친 제3자 \par
ㄹ. 가장의 전세권설정계약에 기하여 등기가 경료된 전세권에 관하여 저당권을 취득한 제3자 \par
\end{framedtext}
\begin{choices}
  \item ㄱ, ㄴ
  \item ㄷ, ㄹ
  \item ㄱ, ㄴ, ㄷ
  \item ㄴ, ㄷ, ㄹ
  \item ㄱ, ㄴ, ㄷ, ㄹ
\end{choices}
\end{question}
\answer{5}
